\documentclass[11pt,a4paper]{article}
\usepackage[margin=1in]{geometry}
\usepackage{amsmath,amssymb}
\usepackage{booktabs}
\usepackage{array}
\usepackage{enumitem}
\usepackage[dvipsnames]{xcolor}
\usepackage{hyperref}
\usepackage{parskip}
\usepackage[T1]{fontenc}
\usepackage{lmodern}
\usepackage{microtype}

\hypersetup{
    colorlinks=true,
    linkcolor=NavyBlue,
    citecolor=NavyBlue,
    urlcolor=NavyBlue,
    pdftitle={Safe Projection onto Linear Safety Constraints - SafeGen-SBP (v2.2.0)},
    pdfauthor={Mobin Khorshidvand}
}

\renewcommand{\familydefault}{\sfdefault}
\emergencystretch=1.5em

\begin{document}

\begin{center}
    {\LARGE \textbf{Safe Projection onto Linear Safety Constraints}}\\[0.5em]
    {\large \textbf{SafeGen-SBP: Engineering Audit, Geometric Foundations, and Numerical Verification (v2.2.0)}}\\[0.75em]
    {\normalsize \textbf{Author:} Mobin Khorshidvand}\\[0.25em]
    {\small \textbf{Final verification environment:} Python 3.12.3 $\cdot$ NumPy 2.1.3 $\cdot$ SciPy 1.14.1}
\end{center}

\vspace{0.5em}
\hrule
\vspace{0.75em}

\section*{Abstract}
SafeGen-SBP (v2.2.0) provides a practical Python implementation for computing Euclidean projections onto linear half-spaces \cite{boyd2004}. The package combines the closed-form single-constraint projection formula with Dykstra's cyclic projection algorithm for intersecting constraints \cite{dykstra1983}, supplemented by an independent feasibility check and benchmark comparisons against SciPy's SLSQP solver. Rather than introducing a novel optimization method or claiming long-term closed-loop safety guarantees, this work focuses on numerical reliability, edge-case auditing, and reproducible verification. All eight unit and integration tests passed during our final verification run. Because our latest benchmarks differed from several figures noted in an earlier draft---specifically regarding raw violation rates and execution speed---this report presents the updated measurements taken directly from the verified execution logs.

\textbf{Keywords:} Euclidean projection; linear constraints; Dykstra's algorithm; numerical verification; reproducibility.

\vspace{0.5em}
\hrule
\vspace{0.75em}

\section{Problem Definition and Geometric Motivation}

In many control and machine learning applications, continuous outputs generated by a model may occasionally violate operational linear inequalities. A straightforward way to enforce these bounds is to apply a post-processing filter that maps any offending vector to the nearest point within the safe set \cite{boyd2004}.

For a single linear inequality, we define the closed half-space as $H = \{x \in \mathbb{R}^n \mid g^T x \le \beta\}$, where $g \neq 0$. The orthogonal (Euclidean) projection of a point $x \in \mathbb{R}^n$ onto $H$ is given in closed form by \cite{boyd2004}:
\[
P_H(x) = x - \frac{\max(0, g^T x - \beta)}{g^T g} g.
\]

Geometrically, if the input vector already satisfies the constraint ($g^T x \le \beta$), the filter leaves it untouched. When a violation occurs, the formula subtracts a scalar multiple of the normal vector $g$, pulling the point orthogonally back onto the boundary hyperplane $g^T x = \beta$. This operator is \emph{piecewise affine}---acting as the identity map inside the feasible region and an affine projection outside---and remains continuous across the boundary, though it is generally non-differentiable where $g^T x = \beta$. Note that this operation stops at the boundary rather than reflecting the vector across it.

\subsection{Two-Dimensional Walkthrough}
To see how this works in practice, consider a simple two-dimensional case with input $x = (4, 5)$, normal vector $g = (1, 1)$, and bound $\beta = 4$. Evaluating the inner product gives $g^T x = 9$, which exceeds $\beta$ by a violation of $9 - 4 = 5$. Since the squared norm of the normal vector is $g^T g = 2$, the projection step scales $g$ by $5 / 2 = 2.5$:
\[
P_H(x) = (4, 5) - 2.5(1, 1) = (1.5, 2.5).
\]

The resulting coordinates sum to $1.5 + 2.5 = 4$, placing the point squarely on the constraint boundary. We include this calculation as Test 1 in our verification suite. Keep in mind that this projection is strictly a pointwise geometric correction at the current step and makes no assumptions about how the underlying system will evolve over subsequent time steps.

\subsection{Scope and Contribution}
Both half-space projection \cite{boyd2004} and Dykstra's algorithm \cite{dykstra1983} are well-established tools in convex analysis, so we make no claims of mathematical novelty here. Instead, the contribution of this work lies in providing a clean, well-documented NumPy implementation, auditing the numerical pitfalls of adding small regularization offsets to the denominator, and supplying a fully reproducible suite of unit tests and solver benchmarks.

\section{Numerical Policy and Multiple Constraints}

\subsection{Squared-Normal Cutoff}
When implementing projection formulas in floating-point arithmetic, nearly zero normal vectors can cause severe numerical instability. To guard against division by underflowed values, our implementation enforces a strict cutoff of $g^T g \le 10^{-14}$ on the \emph{squared} Euclidean norm. We treat this threshold as a practical engineering safeguard rather than a mathematical redefinition of zero.

Under this rule, a normal vector with $\|g\|_2 = 10^{-6}$ yields $g^T g = 10^{-12}$ and is processed normally, whereas a vector with $\|g\|_2 = 10^{-8}$ gives $g^T g = 10^{-16}$ and is rejected as degenerate. Exact zeros, non-finite inputs (\texttt{NaN} or \texttt{Inf}), incompatible dimensions, and invalid tolerances are similarly caught and raised as \texttt{ValueError} exceptions (verified in Test 3 and Test 4). Because multiplying a constraint inequality by a positive scalar changes the magnitude of $g^T g$, users should be mindful of how their constraints are scaled before passing them to the filter.

\subsection{\texorpdfstring{Effect of Adding $10^{-12}$ to the Denominator}{Effect of Adding 1e-12 to the Denominator}}
A common habit in numerical coding is to tack a small constant, such as $\epsilon = 10^{-12}$, onto the denominator $g^T g + \epsilon$ to prevent division by zero. While well-intentioned, this additive offset silently corrupts the projection whenever the constraint normal happens to be small.

For instance, take $g = (10^{-6}, 0)$, $x = (2, 5)$, and $\beta = 10^{-6}$. Here, the squared norm is $g^T g = 10^{-12}$ and the constraint violation is $g^T x - \beta = 10^{-6}$. Using the exact formula, the scaling multiplier is $10^{-6} / 10^{-12} = 10^6$, which shifts $x$ to the true boundary point $(1, 5)$. However, if $10^{-12}$ is added to the denominator, the divisor doubles to $2 \times 10^{-12}$. This cuts the corrective step in half and yields $(1.5, 5)$---a point that still violates the safety constraint.

We verify this exact scenario in Test 4. Note that the helper function \texttt{legacy\_project\_buggy} included in our codebase is simply a pedagogical reconstruction designed to illustrate this specific numerical flaw; it is not meant to imply that any particular third-party library contains this bug.

\subsection{Dykstra's Method and Status Reporting}
When dealing with a polyhedral intersection of multiple constraints $G x \le b$, simply cycling through individual half-space projections (Von Neumann's alternating projections) generally converges to a feasible point, but not necessarily the closest one in Euclidean distance. To find the true metric projection onto the intersection, we implement Dykstra's cyclic projection algorithm \cite{dykstra1983}. Dykstra's method maintains an auxiliary correction vector for each constraint: before projecting onto a given half-space, the algorithm adds that constraint's previous correction back to the current iterate, performs the single-half-space projection, and updates the stored correction as the difference.

In our implementation, we carefully separate convergence criteria---such as step-size tolerances (\texttt{step\_tol}) and iteration limits (\texttt{max\_iters})---from actual constraint satisfaction (\texttt{is\_feasible}). A solution is flagged as certified feasible (\texttt{is\_feasible=True}, returning \texttt{x}) only after an independent post-hoc check confirms that the maximum signed constraint violation across all rows of $G x \le b$ falls at or below \texttt{feas\_tol}.

This distinction is crucial in practice: exhausting the iteration budget (\texttt{budget\_exhausted=True}) does not prove that a problem is infeasible, just as a tiny step between successive iterates does not guarantee that the current point is safe. For example, when fed contradictory constraints ($x_1 + x_2 \le 2$ and $-x_1 - x_2 \le -5$) in Test 7, the solver hits its 50-cycle limit with a maximum signed violation of \texttt{3.0} and properly reports a failure to certify a feasible point (\texttt{x=None}), without falsely claiming a formal certificate of infeasibility.

\section{Experimental Setup and Results}

All empirical results reported below were generated in our final verification run (recorded in \path{results/final_verification.log}, or \path{../results/final_verification.log} relative to \path{docs/}) using Python 3.12.3, NumPy 2.1.3, and SciPy 1.14.1. To ensure consistency, all random sampling was seeded across five fixed runs (\texttt{seeds = [42, 43, 44, 45, 46]}) via \path{numpy.random.default_rng(seed)}.

\subsection{Benchmark 1: Single Half-Space}
In our first benchmark, we tested projection in $\mathbb{R}^{50}$ using a single random constraint direction per seed, normalized to $\|g\|_2 = 10^{-6}$ with bound $\beta = 0$. For each of the five seeds, we drew 1,000 sample vectors from a normal distribution $\mathcal{N}(0.5, 1.0)$, totaling 5,000 test points. Identical inputs were passed to the unfiltered baseline (Raw), the denominator-offset model (Reconstructed $+10^{-12}$), and our exact projection implementation. Table 1 summarizes the mean and population standard deviation of the five per-seed violation rates, alongside the mean of the five per-seed maximum violations.

\begin{table}[h!]
\centering
\caption{Single half-space projection performance across 5,000 samples (5 seeds $\times$ 1,000 inputs).}
\vspace{0.4em}
\small
\begin{tabular}{@{}lcc@{}}
\toprule
\textbf{Method} & \textbf{Violation Rate (Mean $\pm$ SD)} & \textbf{Mean of Per-Seed Max Violation} \\
\midrule
Raw (Unfiltered) & \texttt{52.200000\%} $\pm$ \texttt{20.801250\%} & \texttt{3.334013636005e-06} \\
Reconstructed ($+10^{-12}$ offset) & \texttt{52.200000\%} $\pm$ \texttt{20.801250\%} & \texttt{1.667006818003e-06} \\
Exact projection & \texttt{0.000000\%} $\pm$ \texttt{0.000000\%} & \texttt{1.072201025912e-21} \\
\bottomrule
\end{tabular}
\end{table}

To be fully transparent about how violations were tallied in the code and log: the raw baseline flagged any strictly positive residual ($g^T x - \beta > 0$), the reconstructed-offset test used a threshold of $10^{-15}$, and the exact projection check counted residuals exceeding $10^{-12}$. Because these thresholds differ slightly, they are explicitly disclosed; however, even if evaluated against the strictest threshold ($10^{-15}$), the exact method's worst-case per-seed residual averaged \texttt{1.072201025912e-21}, which is zero up to floating-point rounding.

As seen in Table 1, adding the $10^{-12}$ offset cuts the magnitude of the violation in half (leaving a mean per-seed maximum Euclidean gap of \texttt{1.667006818003} between the offset and exact outputs) but fails to bring the violating points across the boundary. We also note that a raw violation rate of \texttt{74.72\%} $\pm$ \texttt{0.83\%} mentioned in an earlier unverified draft could not be reproduced under this seeding setup (where one random direction $g$ is drawn per seed) and has been discarded in favor of the logged figures above.

\subsection{Benchmark 2: Polytope Projection vs. SLSQP}
Our second benchmark evaluated multi-constraint projection in $\mathbb{R}^{20}$ subject to 21 linear inequalities: an upper bound on the sum of coordinates ($\sum_{i=1}^{20} x_i \le 15$) alongside twenty nonnegativity bounds ($-x_i \le 0$, i.e., $x_i \ge 0$). Across five seeds, we generated 200 points per seed (1,000 paired inputs overall) from $\mathcal{N}(1, 2)$. We configured Dykstra's algorithm with \texttt{step\_tol = 1e-8}, \texttt{feas\_tol = 1e-7}, and \texttt{max\_iters = 500}. For comparison, we solved the equivalent quadratic program $\min_y \frac{1}{2} \|y - x\|_2^2$ subject to $G y \le b$ using SciPy's SLSQP routine with \texttt{ftol = 1e-12} and \texttt{maxiter = 500}. Success counters for Dykstra and SLSQP are tracked independently in the benchmark loop.

Across all 1,000 test instances, Dykstra's method converged to a certified feasible solution every time (\texttt{1000/1000}), requiring an average of \texttt{27.945000} cycles and achieving an overall maximum constraint violation of \texttt{9.999362582391e-08} (within the $10^{-7}$ target tolerance). SciPy's SLSQP solver likewise succeeded on all \texttt{1000/1000} problems with a maximum violation of \texttt{1.044497821567e-12}. Comparing the solutions returned by both solvers across all 1,000 jointly successful cases, the maximum $L_\infty$ difference was only \texttt{1.632513857430e-08}, confirming that both methods arrive at virtually identical projections.

\section{Timing, Discussion, and Limitations}

During the Benchmark 2 run, we recorded wall-clock execution times sequentially using Python's \texttt{time.perf\_counter}, measuring each sample in a single pass without warm-up passes or multi-trial averaging. In our test environment, Dykstra's algorithm averaged \texttt{13.620408 ms} per sample, whereas SciPy's SLSQP took only \texttt{0.475233 ms} per sample---an empirical time ratio (SLSQP / Dykstra) of \texttt{0.034891x}.

This outcome directly contradicts an earlier draft's claim that the pure-NumPy Dykstra loop outperformed SLSQP. In hindsight, the speed gap is unsurprising: SLSQP delegates its heavy lifting to compiled Fortran routines with minimal Python overhead, whereas Dykstra's cyclic updates execute hundreds of small array operations inside a Python loop for each cycle. Of course, single-pass timings are sensitive to hardware, BLAS threading, and library versions, so these numbers should be read as an environment-specific observation rather than a universal benchmark.

\subsection{Verification Summary}
All eight automated tests in our verification suite passed without issue:
\begin{enumerate}[leftmargin=2em,itemsep=0.25em]
    \item \textbf{Test 1:} The 2D analytical walkthrough (\texttt{[4, 5] -> [1.5, 2.5]} on $x_1 + x_2 \le 4$) from Section 1.1.
    \item \textbf{Test 2:} Bitwise preservation of already-feasible inputs (\texttt{[1, 2]} on $x_1 + x_2 \le 4$).
    \item \textbf{Test 3:} Explicit rejection (\texttt{ValueError}) of zero, \texttt{NaN}, and \texttt{Inf} normal vectors, incompatible matrix/vector dimensions, and invalid tolerances or iteration budgets.
    \item \textbf{Test 4:} Verification of the $10^{-14}$ squared-norm cutoff (accepting $\|g\|_2 = 10^{-6}$, rejecting $\|g\|_2 = 10^{-8}$) and the reconstructed $+10^{-12}$ denominator offset flaw (\texttt{[2, 5] -> [1.5, 5]} instead of \texttt{[1, 5]}).
    \item \textbf{Test 5:} Large-scale 2D normal projection ($x = [4, 3]$ with $g = [10^6, 0]$ and $\beta = 2 \times 10^6$ projecting accurately to \texttt{[2, 3]}).
    \item \textbf{Test 6:} Tight agreement with SciPy's SLSQP solver: when running Dykstra with \texttt{step\_tol = 1e-12} and \texttt{feas\_tol = 1e-12}, the $L_\infty$ discrepancy against SLSQP measured \texttt{1.152855588771e-12} (well below the explicit \texttt{< 1.1e-9} assertion).
    \item \textbf{Test 7:} Graceful handling of contradictory constraints ($x_1 + x_2 \le 2$ and $-x_1 - x_2 \le -5$, yielding signed max violation \texttt{3.0} after 50 cycles) without false feasibility or infeasibility-certificate claims.
    \item \textbf{Test 8:} Explicit separation of iteration-budget exhaustion (\texttt{max\_iters = 1}, \texttt{converged = False}, \texttt{is\_feasible = False}, \texttt{max\_violation = 2.32}) from safety certification.
\end{enumerate}

\subsection{Limitations}
To keep the scope of this work clear, several practical limitations should be noted. All experiments were conducted in standard \texttt{float64} precision using synthetic Gaussian vectors and static linear inequalities. We have not evaluated this filter in hardware-in-the-loop setups, real robotic trajectory tracking, nonlinear obstacle avoidance, or closed-loop dynamical systems.

Furthermore, projecting a single output onto a convex polytope only guarantees constraint satisfaction at that instant; in a dynamical system, one-step feasibility does not prevent the system from entering a state where future constraints become impossible to satisfy (lack of recursive feasibility). Finally, as noted in Section 2.3, hitting the maximum iteration cap in Dykstra's algorithm is simply a timeout condition, not a formal proof that the constraint set is empty. Complete execution logs and reproduction scripts are available in the accompanying repository materials.

\begin{thebibliography}{9}
\bibitem{boyd2004}
Boyd, S., \& Vandenberghe, L. (2004). \emph{Convex Optimization}. Cambridge University Press.

\bibitem{dykstra1983}
Dykstra, R. L. (1983). An Algorithm for Restricted Least Squares Regression. \emph{Journal of the American Statistical Association}, 78(384), 837--842.
\end{thebibliography}

\end{document}
